%%=============================================================================
%% Requirementanalyse
%%=============================================================================

\chapter{\IfLanguageName{dutch}{Requirements-analyse}{Requirements analysis}}%
\label{ch:requirementsanalyse}

De voorgaande literatuurstudie bood inzicht in de theoretische achtergrond van burn-out en bracht de bestaande applicaties in kaart die in wetenschappelijk gescreende databanken gedocumenteerd zijn. Hoewel deze inventarisatie aantoont welke oplossingen reeds onderzocht zijn, ontbrak de vertaling naar specifieke, klinisch onderbouwde functionaliteiten voor de beoogde doelgroep van Gen Z. Om de brug te slaan tussen deze theoretische verkenning en de daadwerkelijke requirementsanalyse, werd een interview afgenomen met expert en co-promotor Mevr. K. Van den Broeck (klinisch psycholoog).

\vspace{1em}

Het doel van dit interview was tweeledig:

\begin{enumerate}
    \item \textbf{Verdieping van het probleemdomein}: Het verkrijgen van nuances uit de klinische praktijk die specifiek relevant zijn voor de doelgroep, zoals het onderscheid tussen chronische overspanning en burn-out.
    \item \textbf{Verzamelen van requirements}: Het identificeren van essentiële (niet-) \newline functionele vereisten die de basis vormen voor de later opgestelde \newline ontwerpprincipes.
\end{enumerate}

Het gesprek vond plaats via een semi-gestructureerd interviewformaat om ruimte te bieden voor informatie vanuit de praktijkervaring. De gehanteerde vragenlijst is terug te vinden in bijlage \ref{ch:interview-requirementsanalyse}. De voornaamste bevindingen die fungeren als directe input voor de requirements worden hieronder nader toegelicht.

\newpage

\section{Kerninzichten voor de requirements}
\label{sec:inzichten-voor-requirements}

\vspace{1em}

De synthese van dit interview leidde tot de volgende cruciale pijlers die de verdere bepaling van de functionele en niet-functionele requirements hebben gestuurd:

\begin{itemize}
    \item \textbf{Focus op de 'externe rem':} De doelgroep bevindt zich vaak in een staat van chronische overprikkeling (hyper-alertheid). De applicatie moet fungeren als een preventieve tool die helpt bij het herkennen van fysiologische signalen voordat het stadium van burn-out bereikt wordt.
    \item \textbf{Methodiek via het 5G-model:} Voor de zelfreflectie wordt de cognitieve gedragstherapeutische methode van het 5G-model (Gebeurtenis, Gedachte, Gevoel, Gedrag, Gevolg) gehanteerd. Dit wordt vertaald naar een geleide dagboekfunctie die de gebruiker helpt om dieperliggende patronen zoals perfectionisme te ontleden.
    \item \textbf{Fase-gebaseerde content (Chunking):} De interventie moet stapsgewijs verlopen om information overload te vermijden. Dit kan bijvoorbeeld voor secondaire burn-out preventie door eerst de focus te leggen op 'leren remmen' (ademhaling, zintuiglijke oefeningen) en pas bij voldoende fysieke rust schakelt de app over naar cognitieve reflectie over de oorzaken. Echter wordt de focus gelegd op secondaire preventie van chronische overspanning en zal dit een andere vorm kunnen krijgen.
    \item \textbf{Eliminatie van prestatiedruk:} Een essentieel inzicht is het vermijden van \textit{gamification} die stress kan opwekken. De app mag geen gebruik maken van scores, streaks, ratings of rankings. Deze elementen activeren het stresssysteem bij perfectionistische gebruikers en kunnen leiden tot een gevoel van falen wanneer de app niet 'perfect' gebruikt wordt.
    \item \textbf{Risico's van sociale ondersteuning:} Hoewel verbondenheid waardevol is, moet er opgepast worden voor directe sociale vergelijking (bijv. "Waarom herstelt de ander sneller?"). Sociale functies moeten daarom anoniem en niet-competitief zijn, bijvoorbeeld door middel van het delen van algemene, anonieme bemoedigende boodschappen in plaats van directe chatfuncties.
    \item \textbf{Retentie door psycho-educatie:} Om vroegtijdig afhaken te voorkomen (het zoeken naar een \textit{quick fix}), is psycho-educatie cruciaal. Door uit te leggen hoe het zenuwstelsel werkt, begrijpt de gebruiker waarom resultaten uitblijven als er geen rust wordt genomen, wat de motivatie op lange termijn verhoogt.
    \item \textbf{Personalisatie van notificaties:} Om te vermijden dat de app als een extra 'verplichting' aanvoelt, moeten notificaties instelbaar zijn op basis van de individuele behoefte en de mate van smartphonegebruik van de gebruiker.
\end{itemize}

De bovenstaande inzichten, samen met de kennis die voortkwam uit de literatuurstudie, vormen de directe onderbouwing voor de geprioriteerde requirements die in de volgende sectie (\ref{sec:requirements}) worden opgesomd.

\vspace{1em}

\section{Requirements}
\label{sec:requirements}

\subsection{Functionele requirements}
\label{subsec:frs}

\vspace{1em}

De functionele requirements (\acrshort{fr}'s) beschrijven de specifieke acties, gedragingen en functies die de applicatie moet kunnen uitvoeren. Ze definiëren met andere woorden de directe interactie tussen de gebruiker en het systeem en bepalen wat de applicatie precies 'doet' om de doelstelling (verminderen van chronische overspanning) te faciliteren.

\vspace{1em}

\subsubsection{User stories}
\label{subsubsec:user-stories}

\vspace{1em}

Om deze requirements te identificeren werd gebruikgemaakt van user stories als ondersteunende techniek. Deze user stories zijn opgesteld met een focus op de preventie van chronische overspanning en het ondersteunen van stressregulatie, eerder dan op de behandeling van burn-out. Op basis van deze user stories worden de uiteindelijke functionele requirements, aangevuld met een MoSCoW-prioritering, samengebracht in een overzichtelijke tabel. Daarnaast kunnen bepaalde user stories, in combinatie met inzichten uit de literatuurstudie, later bijdragen aan het formuleren van ontwerpprincipes.

\vspace{1em}

\paragraph{Gebruiker}

\vspace{1em}

\textit{Account en basisgebruik}
\hrule
\vspace{1em}

Als \textbf{gebruiker}, \newline
wil ik \textbf{mij kunnen registreren en aanmelden}, \newline
zodat \textit{ik toegang krijg tot mijn persoonlijke traject}.

\vspace{1em}

Als \textbf{gebruiker}, \newline
wil ik \textbf{juridische informatie en het privacy beleid kunnen raadplegen}, \newline
zodat \textit{ik weet hoe mijn gegevens worden gebruikt}.

\vspace{1em}
\vspace{1em}

\textit{Onboarding en personalisatie}
\hrule
\vspace{1em}

Als \textbf{gebruiker}, \newline
wil ik \textbf{een korte en visuele onboarding}, \newline
zodat \textit{ik snel kan starten zonder overweldigd te worden}.

\vspace{1em}

Als \textbf{gebruiker}, \newline
wil ik \textbf{mijn profiel en dashboard kunnen personaliseren}, \newline
zodat \textit{de app aansluit bij mijn voorkeuren en noden}.

\vspace{1em}

Als \textbf{gebruiker}, \newline
wil ik \textbf{kunnen kiezen hoe intensief ik begeleid word}, \newline
zodat \textit{ik autonomie kan behouden in mijn traject}.

\vspace{1em}
\vspace{1em}

\textit{Zelfreflectie en dagboek}
\hrule
\vspace{1em}

Als \textbf{gebruiker}, \newline
wil ik \textbf{mijn stressniveau kunnen instellen}, \newline
zodat \textit{ik mijn stressniveaus en trends kan opvolgen}.

\vspace{1em}

Als \textbf{gebruiker}, \newline
wil ik \textbf{mijn gedachten/situaties kunnen registreren in een dagboek}, \newline
zodat \textit{ik inzicht krijg in mijn gedrag en emoties}.

\vspace{1em}

Als \textbf{gebruiker}, \newline
wil ik \textbf{mijn eerdere dagboeknotities kunnen raadplegen}, \newline
zodat \textit{ik evoluties en patronen kan herkennen}.

\vspace{1em}

Als \textbf{gebruiker}, \newline
wil ik \textbf{mijn dagboekgegevens kunnen exporteren}, \newline
zodat \textit{ik deze kan delen met een hulpverlener}.

\vspace{1em}

Als \textbf{gebruiker}, \newline
wil ik \textbf{dat mijn ingevoerde gegevens worden samengevat of geanalyseerd}, \newline
zodat \textit{ik terugkerende patronen kan herkennen}.

\vspace{1em}
\vspace{1em}

\textit{Psycho-educatie}
\hrule
\vspace{1em}

Als \textbf{gebruiker}, \newline
wil ik \textbf{toegang tot korte en begrijpbare educatieve content}, \newline
zodat \textit{ik inzicht krijg in stress, coping en chronische overspanning}.

\vspace{1em}

Als \textbf{gebruiker}, \newline
wil ik \textbf{content kunnen raadplegen op mijn eigen tempo}, \newline
zodat \textit{ik informatie beter kan verwerken}.

\vspace{1em}

Als \textbf{gebruiker}, \newline
wil ik \textbf{op relevante momenten educatieve feedback ontvangen}, \newline
zodat \textit{ik realistische verwachtingen heb over mijn traject}.

\vspace{1em}
\vspace{1em}

\textit{Oefeningen en stressreductie}
\hrule
\vspace{1em}

Als \textbf{gebruiker}, \newline
wil ik \textbf{toegang tot korte oefeningen} (zoals ademhaling of mindfulness), \newline
zodat \textit{ik mijn stress onmiddellijk kan verlagen}.

\vspace{1em}

Als \textbf{gebruiker}, \newline
wil ik \textbf{oefeningen in kleine stappen kunnen doorlopen} (chunking), \newline
zodat \textit{ik me competent voel en niet overweldigd raak}.

\vspace{1em}

Als \textbf{gebruiker}, \newline
wil ik \textbf{kunnen kiezen uit oefeningen die passen binnen de fase van mijn traject}, \newline
zodat \textit{ze relevant zijn voor mijn traject}.

\vspace{1em}

Als \textbf{gebruiker}, \newline
wil ik \textbf{begeleid worden via reflectievragen rond lichamelijke signalen}, \newline
zodat \textit{ik mijn grenzen beter kan herkennen}.

\vspace{1em}

Als \textbf{gebruiker}, \newline
wil ik \textbf{begeleide rustmomenten of pauzes aangeboden krijgen}, \newline
zodat \textit{mijn stresssysteem kan herstellen}.

\vspace{1em}

Als \textbf{gebruiker}, \newline
wil ik \textbf{een focusmodus kunnen activeren}, \newline
zodat \textit{ik tijdelijk minder afleiding ervaar}.

\vspace{1em}

Als \textbf{gebruiker}, \newline
wil ik \textbf{toegang tot een noodfunctie bij acute stress}, \newline
zodat \textit{ik snel ondersteuning krijg in moeilijke momenten}.

\vspace{1em}
\vspace{1em}

\textit{Feedback en progressie}
\hrule
\vspace{1em}

Als \textbf{gebruiker}, \newline
wil ik \textbf{mijn voortgang visueel kunnen volgen}, \newline
zodat \textit{ik gemotiveerd blijf}.

\vspace{1em}

Als \textbf{gebruiker}, \newline
wil ik \textbf{positieve feedback krijgen op mijn acties}, \newline
zodat \textit{ik gemotiveerd blijf en mijn doelen voor ogen houd}.

\vspace{1em}

Als \textbf{gebruiker}, \newline
wil ik \textbf{kleine beloningen ontvangen} (zoals badges of mijlpalen), \newline
zodat \textit{mijn vooruitgang zichtbaar wordt zonder prestatiedruk}.

\vspace{1em}
\vspace{1em}

\textit{Notificaties en nudging}
\hrule
\vspace{1em}

Als \textbf{gebruiker}, \newline
wil ik \textbf{gepersonaliseerde notificaties ontvangen}, \newline
zodat \textit{ze relevant zijn voor mijn situatie}.

\vspace{1em}

Als \textbf{gebruiker}, \newline
wil ik \textbf{controle hebben over frequentie en timing van notificaties}, \newline
zodat \textit{ik niet overweldigd raak}.

\vspace{1em}

Als \textbf{gebruiker}, \newline
wil ik \textbf{interactieve notificaties ontvangen} (vb. quick check-in), \newline
zodat \textit{ik betrokken blijf in mijn traject}.

\vspace{1em}

Als \textbf{gebruiker}, \newline
wil ik \textbf{nudges krijgen die afgestemd zijn op mijn gedrag}, \newline
zodat \textit{ik gemotiveerd blijf zonder druk te voelen}.

\vspace{1em}
\vspace{1em}

\textit{Verbondenheid en community}
\hrule
\vspace{1em}

Als \textbf{gebruiker}, \newline
wil ik \textbf{op een laagdrempelige manier in contact kunnen komen met anderen}, \newline
zodat \textit{ik me minder alleen voel en sociale steun ervaar}.

\vspace{1em}

Als \textbf{gebruiker}, \newline
wil ik \textbf{dat sociale interacties geen prestaties of vergelijkingen tonen}, \newline
zodat \textit{prestatiedruk vermeden wordt}.

\vspace{1em}

Als \textbf{gebruiker}, \newline
wil ik \textbf{anonieme aanmoedigingen kunnen versturen en ontvangen}, \newline
zodat \textit{ik steun kan geven en krijgen op een laagdrempelige manier en zo sociale steun ervaar}.

\vspace{1em}

Als \textbf{gebruiker}, \newline
wil ik \textbf{kunnen bepalen hoeveel informatie ik deel met anderen}, \newline
zodat \textit{mijn privacy gewaarborgd blijft}.

\vspace{1em}

Als \textbf{gebruiker}, \newline
wil ik \textbf{anoniem en zonder zichtbare prestatie-indicatoren kunnen interageren}, \newline
zodat \textit{sociale vergelijking vermeden wordt}.

\vspace{1em}
\vspace{1em}

\textit{Motivatie en engagement}
\hrule
\vspace{1em}

Als \textbf{gebruiker}, \newline
wil ik \textbf{persoonlijke doelen kunnen instellen} (SMARTER), \newline
zodat \textit{ik een traject op maat kan volgen en gericht kan werken aan verbetering}.

\vspace{1em}

Als \textbf{gebruiker}, \newline
wil ik \textbf{een progressiesysteem of verhaallijn ervaren}, \newline
zodat \textit{de applicatie boeiend blijft}.

\vspace{1em}
\vspace{1em}

\textit{Hulp en doorverwijzing}
\hrule
\vspace{1em}

Als \textbf{gebruiker}, \newline
wil ik \textbf{toegang tot informatie over professionele hulpverleners}, \newline
zodat \textit{ik weet waar ik terechtkan}.

\vspace{1em}

Als \textbf{gebruiker}, \newline
wil ik \textbf{advies krijgen om professionele hulp te zoeken wanneer nodig}, \newline
zodat \textit{ik tijdig geholpen word bij verergering}.

\vspace{1em}

\paragraph{Admin}

\vspace{1em}

Als \textbf{admin}, \newline
wil ik \textbf{gebruikers kunnen beheren}, \newline
zodat \textit{de veiligheid van het platform gewaarborgd blijft}.

\vspace{1em}

Als \textbf{admin}, \newline
wil ik \textbf{community-interacties kunnen modereren}, \newline
zodat \textit{de omgeving respectvol en veilig blijft}.

\vspace{1em}

Als \textbf{admin}, \newline
wil ik \textbf{informatie over hulpinstanties kunnen beheren}, \newline
zodat \textit{gebruikers correcte en actuele info krijgen}.

\vspace{1em}

\paragraph{Expert}

\vspace{1em}

Als \textbf{expert}, \newline
wil ik \textbf{mij kunnen registreren als professionele gebruiker}, \newline
zodat \textit{ik content kan bijdragen}.

\vspace{1em}

Als \textbf{expert}, \newline
wil ik \textbf{educatieve content kunnen toevoegen en aanpassen}, \newline
zodat \textit{gebruikers betrouwbare informatie krijgen}.

\vspace{1em}

Als \textbf{expert}, \newline
wil ik \textbf{content kunnen valideren}, \newline
zodat \textit{de kwaliteit van de informatie gewaarborgd blijft}.

\vspace{1em}
\vspace{1em}
\vspace{1em}

\subsubsection{Functionele requirements met MoSCoW}
\label{subsubsec:frs-moscow}

\vspace{2em}

{\small
    \textcolor{green!80!black}{\textbf{M}}ust have \quad
    \textcolor{orange}{\textbf{S}}hould have \quad
    \textcolor{blue}{\textbf{C}}ould have \quad
    \textcolor{black!60}{\textbf{W}}on’t have
}

{
    \small
    \setlength\tabcolsep{3pt}
    \begin{longtable}{
            >{\raggedright\arraybackslash}p{3.0cm}
            >{\raggedright\arraybackslash}p{10.5cm}
            >{\centering\arraybackslash}p{2.0cm}
        }
        \toprule
        \textbf{Rol} & \textbf{Functionaliteit} & \textbf{Prioriteit} \\
        \midrule
        \endfirsthead
        
        \multicolumn{3}{c}{\tablename\ \thetable\ -- vervolg van vorige pagina} \\
        \toprule
        \textbf{Rol} & \textbf{Functionaliteit} & \textbf{Prioriteit} \\
        \midrule
        \endhead
        
        \midrule \multicolumn{3}{r}{Vervolg op volgende pagina} \\
        \endfoot
        
        \bottomrule
        \caption[Overzicht requirements]{Oplijsting van functionele requirements met MoSCoW-prioriteit.}
        \label{tab:requirements}
        \endlastfoot
        \addlinespace[6pt]
        
        %==============================
        Gebruiker & Registreren en aanmelden & \textcolor{green!80!black}{\textbf{M}} \\
        Alle rollen & Raadplegen privacybeleid en juridische info & \textcolor{green!80!black}{\textbf{M}} \\
        Admin & Beheren van gebruikers & \textcolor{green!80!black}{\textbf{M}} \\
        Admin & Modereren van community & \textcolor{green!80!black}{\textbf{M}} \\
        Admin & Beheren van hulpinstanties & \textcolor{green!80!black}{\textbf{M}} \\
        Expert & Registreren als expert & \textcolor{green!80!black}{\textbf{M}} \\
        Expert & Toevoegen van educatieve content & \textcolor{blue}{\textbf{C}} \\
        Expert & Aanpassen van educatieve content & \textcolor{orange}{\textbf{S}} \\
        Expert & Valideren van educatieve content & \textcolor{orange}{\textbf{S}} \\ \addlinespace[20pt]
        
        \multicolumn{3}{l}{\textbf{Onboarding \& personalisatie}} \\ \addlinespace[6pt]
        Gebruiker & Doorlopen van visuele onboarding & \textcolor{green!80!black}{\textbf{M}} \\
        Gebruiker & Personaliseren van profiel & \textcolor{blue}{\textbf{C}} \\
        Gebruiker & Personaliseren van dashboard & \textcolor{orange}{\textbf{S}} \\
        Gebruiker & Instellen van begeleidingsintensiteit & \textcolor{orange}{\textbf{S}} \\
        Gebruiker & Instellen van persoonlijke doelen & \textcolor{green!80!black}{\textbf{M}} \\
        Gebruiker & Instellen van notificatiefrequentie & \textcolor{green!80!black}{\textbf{M}} \\
        Gebruiker & Personaliseren van het aanbod aan oefeningen a.d.h.v. vorderingen binnen het persoonlijk traject & \textcolor{orange}{\textbf{S}} \\ \addlinespace[20pt]
        
        \multicolumn{3}{l}{\textbf{Zelfreflectie \& dagboek}} \\ \addlinespace[6pt]
        Gebruiker & Registreren van stressniveau & \textcolor{green!80!black}{\textbf{M}} \\
        Gebruiker & Bijhouden van gestructureerd dagboek & \textcolor{green!80!black}{\textbf{M}} \\
        Gebruiker & Raadplegen van dagboekhistoriek & \textcolor{green!80!black}{\textbf{M}} \\
        Gebruiker & Exporteren van dagboekgegevens & \textcolor{orange}{\textbf{S}} \\
        Gebruiker & Automatische analyse / samenvatting van gegevens & \textcolor{orange}{\textbf{S}} \\ \addlinespace[20pt]
        
        \multicolumn{3}{l}{\textbf{Psycho-educatie}} \\ \addlinespace[6pt]
        Gebruiker & Raadplegen van educatieve content & \textcolor{green!80!black}{\textbf{M}} \\
        Gebruiker & Ontvangen van contextuele educatieve feedback & \textcolor{orange}{\textbf{S}} \\ \addlinespace[20pt]
        
        \multicolumn{3}{l}{\textbf{Oefeningen \& stressregulatie}} \\ \addlinespace[6pt]
        Gebruiker & Raadplegen van stressreducerende oefeningen & \textcolor{green!80!black}{\textbf{M}} \\
        Gebruiker & Raadplegen van begeleide rustmomenten & \textcolor{orange}{\textbf{S}} \\ \addlinespace[20pt]
        
        \multicolumn{3}{l}{\textbf{Feedback \& progressie}} \\ \addlinespace[6pt]
        Gebruiker & Raadplegen van een visuele weergave van de huidige trajectvoortgang & \textcolor{green!80!black}{\textbf{M}} \\
        Gebruiker & Ontvangen van positieve/informatieve feedback & \textcolor{green!80!black}{\textbf{M}} \\
        Gebruiker & Niet-competitieve beloningen & \textcolor{blue}{\textbf{C}} \\ \addlinespace[20pt]
        
        \multicolumn{3}{l}{\textbf{Notificaties \& nudging}} \\ \addlinespace[6pt]
        Gebruiker & Ontvangen van gepersonaliseerde notificaties & \textcolor{orange}{\textbf{S}} \\
        Gebruiker & Interactieve notificaties & \textcolor{blue}{\textbf{C}} \\
        Gebruiker & Gedragsgerichte nudges ontvangen & \textcolor{orange}{\textbf{S}} \\ \addlinespace[20pt]
        
        \multicolumn{3}{l}{\textbf{Verbondenheid \& community}} \\ \addlinespace[6pt]
        Gebruiker & Interageren op een laagdrempelige manier zonder prestatie- of vergelijkingsmechanismen & \textcolor{blue}{\textbf{C}} \\
        Gebruiker & Versturen/ontvangen van anonieme aanmoedigingen & \textcolor{orange}{\textbf{S}} \\
        Gebruiker & Controle over zichtbaarheid en gedeelde informatie & \textcolor{green!80!black}{\textbf{M}} \\ \addlinespace[20pt]
        
        \multicolumn{3}{l}{\textbf{Motivatie \& engagement}} \\ \addlinespace[6pt]
        Gebruiker & Progressiesysteem / verhaallijn & \textcolor{orange}{\textbf{S}} \\ \addlinespace[20pt]
        
        \multicolumn{3}{l}{\textbf{Hulp \& doorverwijzing}} \\ \addlinespace[6pt]
        Gebruiker & Raadplegen van professionele hulpopties & \textcolor{green!80!black}{\textbf{M}} \\
        Gebruiker & Ontvangen van advies tot doorverwijzing & \textcolor{green!80!black}{\textbf{M}} \\ \addlinespace[20pt]
        
    \end{longtable}
}

\subsection{Niet-functionele requirements}
\label{subsec:nfrs}

\vspace{1em}

Na het definiëren van de functionele requirements, die beschrijven wat de applicatie moet doen, wordt in deze sectie gefocust op de niet-functionele requirements (\acrshort{nfr}’s). Deze specificeren hoe het systeem moet functioneren en onder welke kwaliteitsvoorwaarden de functionaliteiten worden aangeboden. In de context van de preventie van chronische overspanning zijn deze cruciaal, omdat de manier waarop informatie wordt aangeboden direct invloed heeft op het stressniveau van de gebruiker.

\vspace{1em}

De \acrshort{nfr}’s worden opgesteld aan de hand van het ISO/IEC 25010-kwaliteitsmodel \parencite{iso25010_2023} en worden geformuleerd volgens het SMART-principe.

\vspace{1em}
\vspace{1em}
\hrule
\vspace{1em}

\subsubsection{Functionele geschiktheid}
\label{subsubsec:functionele-geschiktheid}

\vspace{1em}

\paragraph{Functionele compleetheid}
\vspace{1em}

\begin{itemize}
    \item Alle kernfunctionaliteiten \textit{(dagboek, stressregistratie, oefeningen, psycho- \newline educatie, voortgang)} zijn onmiddellijk beschikbaar zonder verplicht gebruik van community-functionaliteiten of betaalmuren.
\end{itemize}

\paragraph{Functionele correctheid}
\vspace{1em}

\begin{itemize}
    \item Berekeningen van stressniveaus en trends moeten in 99\% van de gevallen consistent en reproduceerbaar zijn bij identieke invoer.
    \item 100\% van feedback en aanbevelingen moet gebaseerd zijn op actuele en gevalideerde gebruikersdata.
    \item De weergegeven informatie van beschikbare psychologen moet in 95\% van de gevallen correct zijn.
    \item 100\% van psycho-educatieve content moet vooraf gevalideerd (of geschreven) zijn door erkende experts vóór publicatie.
\end{itemize}

\vspace{1em}
\vspace{1em}
\hrule
\vspace{1em}

\subsubsection{Prestatie-efficiëntie}
\label{subsubsec:prestatie-efficientie}

\vspace{1em}

\paragraph{Snelheid}
\vspace{1em}

\begin{itemize}
    \item 95\% van alle schermen moet laden binnen 1,5 seconde, en 100\% binnen 3 \newline seconden bij normaal netwerkgebruik.
\end{itemize}

\vspace{1em}
\vspace{1em}
\vspace{1em}

\subsubsection{Uitwisselbaarheid}
\label{subsubsec:uitwisselbaarheid}

\vspace{1em}

\paragraph{Koppelbaarheid}
\vspace{1em}

\begin{itemize}
    \item Dagboek data moet exporteerbaar zijn in PDF en CSV-formaat.
    \item Exportbestanden mogen maximaal 0,01\% dataverlies bevatten \textit{($\leq$ 1 fout per 10.000 records)}.
\end{itemize}

\vspace{1em}
\vspace{1em}
\hrule
\vspace{1em}

\subsubsection{Bruikbaarheid}
\label{subsubsec:bruikbaarheid}

\vspace{1em}

\paragraph{Herkenbare geschiktheid}
\vspace{1em}

\begin{itemize}
    \item 95\% van nieuwe gebruikers moet binnen 10 seconden correct kunnen aangeven wat de primaire functie van de app is.
\end{itemize}

\paragraph{Leerbaarheid}
\vspace{1em}

\begin{itemize}
    \item 95\% van nieuwe gebruikers voltooit onboarding binnen 5 minuten zonder externe hulp.
    \item 95\% van gebruikers start binnen 10 minuten zelfstandig een eerste oefening.
    \item 90\% van de gebruikers moet, na het doornemen van een psycho-educatief artikel, 80\% van de content in eigen woorden kunnen uitleggen zonder het raadplegen van een zoekmachine.
\end{itemize}

\paragraph{Bedienbaarheid}
\vspace{1em}

\begin{itemize}
    \item 95\% van gebruikers kan binnen 5 minuten een dagboekentry aanmaken, ook bij eerste gebruik.
    \item 95\% van gebruikers kan binnen 1 minuut stressniveau registreren, ook bij eerste gebruik.
    \item 95\% van gebruikers kan binnen 1 minuut een oefening starten zonder hulp, ook bij eerste gebruik.
\end{itemize}

\paragraph{Toegankelijkheid}
\vspace{1em}

\begin{itemize}
    \item De app voldoet volledig aan WCAG 2.1 niveau AA, geverifieerd via een accessibility checker.
\end{itemize}

\vspace{1em}
\vspace{1em}
\hrule
\vspace{1em}

\subsubsection{Beveiligbaarheid}
\label{subsubsec:beveiligbaarheid}

\vspace{1em}

\paragraph{Vertrouwelijkheid}
\vspace{1em}

\begin{itemize}
    \item 100\% van gevoelige gegevens (dagboek, stressdata, voortgang) is enkel toegankelijk na succesvolle authenticatie (login vereist).
    \item 100\% van psycho-educatieve content toegevoegd door experts wordt pas zichtbaar na admin-goedkeuring.
\end{itemize}

\paragraph{Integriteit}
\vspace{1em}

\begin{itemize}
    \item Het systeem kan aantonen dat een dagboekentry effectief werd ingevoerd door de ingelogde gebruiker.
\end{itemize}

\paragraph{Authenticiteit}
\vspace{1em}

\begin{itemize}
    \item Bij het registreren van een expert dienen de voornaam, achternaam en contactgegevens gekend te zijn.
\end{itemize}

\vspace{1em}
\vspace{1em}
\hrule
\vspace{1em}

\subsubsection{Onderhoudbaarheid}
\label{subsubsec:onderhoudbaarheid}

\vspace{1em}

\paragraph{Herbruikbaarheid}
\vspace{1em}

\begin{itemize}
    \item 100\% van de UI- en UX-uitwerking van de applicatie moet gebaseerd zijn op een vastgelegd en herbruikbaar ontwerpkader (ontwerpprincipes), zodat UI-componenten en interactiepatronen consistent kunnen worden hergebruikt binnen andere schermen of toekomstige m-healthtoepassingen.
\end{itemize}

\paragraph{Wijzigbaarheid}
\vspace{1em}

\begin{itemize}
    \item 95\% van psycho-educatieve content moet door geautoriseerde experts kunnen worden aangepast binnen 15 minuten via een WYSIWYG-interface, en dus zonder codewijziging.
    \item Gebruikers moeten op elk moment de frequentie van notificaties kunnen wijzigen via instellingen, zonder dat hiervoor een update of herconfiguratie van het systeem nodig is.
    \item De startpagina of het dashboard moet zodanig modulair worden opgebouwd dat gebruikers functionaliteiten kunnen tonen, verbergen of herschikken naar persoonlijke voorkeur.
\end{itemize}

\vspace{1em}
\vspace{1em}
\hrule
\vspace{1em}

\subsubsection{Betrouwbaarheid}
\label{subsubsec:betrouwbaarheid}

\vspace{1em}

\paragraph{Beschikbaarheid}
\vspace{1em}

\begin{itemize}
    \item De volgende kernfunctionaliteiten moeten in minstens 80\% van de testscenario’s offline beschikbaar zijn via lokale opslag en caching: \textit{dagboek, oefeningen en stressregistratie}. Gegevens worden gesynchroniseerd zodra opnieuw een internetverbinding beschikbaar is.
\end{itemize}

\begin{itemize}
    \item Kernfunctionaliteiten \textit{(dagboek, oefeningen, stressregistratie)} moeten 100\% offline beschikbaar zijn.
\end{itemize}

\paragraph{Foutbestendigheid}
\vspace{1em}

\begin{itemize}
    \item Uit gebruikerscomfort en stressreductie mag het systeem nooit crashen door foutieve gebruikersinput, netwerkonderbreking, enz.
\end{itemize}

\vspace{1em}
\vspace{1em}
\hrule
\vspace{1em}

\subsubsection{Overdraagbaarheid}
\label{subsubsec:overdraagbaarheid}

\vspace{1em}

\paragraph{Installeerbaarheid}
\vspace{1em}

\begin{itemize}
    \item De applicatie moet compatibel zijn met iOS en Android.
    \item 95\% van de gebruikers kan de app binnen 1 minuut eenvoudig vinden en \newline identificeren als relevant voor hun noden, zonder externe hulp of \newline bijkomende uitleg.
    \item De app moet binnen 3 minuten geïnstalleerd en voor het eerst succesvol geopend kunnen worden.
\end{itemize}

\newpage




\section{Ontwerpprincipes}
\label{sec:ontwerpprincipes}

\vspace{1em}

De literatuurstudie toont aan dat de een \acrshort{mhealth} app moet inzetten op stressreducerend ontwerp dat de mentale rust bevordert en de gebruikers ondersteunt bij het vinden van een \textit{off switch}. Een lage cognitieve belasting vormt hierbij een essentieel uitgangspunt. 

\vspace{1em}

Op basis van voorgaand onderzoek werden ontwerpprincipes opgesteld voor zowel deze \acrshort{poc} als het bredere ontwerp van prikkelarme en mentaal gezonde digitale toepassingen. Deze principes zijn gericht op het verminderen van cognitieve belasting, het bevorderen van mentale rust en het voorkomen van onnodige stress tijdens het gebruik van de app. Om een gestructureerd overzicht te bieden, werden de principes onderverdeeld in vier categorieën:

\vspace{1em}

\begin{itemize}[nosep]
    \item \textbf{Development}: technisch vertrouwen en het vermijden van technostress.
    \item \textbf{Design}: eenvoud en een lage cognitieve belasting.
    \item \textbf{Interaction}: motivatietechnieken en ondersteuning van \acrshort{sdt}.
    \item \textbf{Content}: relevante en ondersteunende inhoud.
\end{itemize}

\vspace{1em}

\subsubsection{Development}
\label{subsubsec:principes-development}

{
    \small
    \setlength\tabcolsep{3pt}
    \begin{longtable}{
            >{\raggedright\arraybackslash}p{0.2cm}  % Nummer
            >{\raggedright\arraybackslash}p{3.3cm}  % Principe (korte naam regel)
            @{\hspace{0.5cm}}
            >{\raggedright\arraybackslash}p{8.4cm}  % Toepassing (hoe toe te passen)
            @{\hspace{0.5cm}}
            >{\raggedright\arraybackslash}p{2.7cm}  % Theoretische basis (waarop werd dit gebaseerd?)
        }
        
        \toprule
        \multicolumn{2}{l}{\textbf{Principe}} & \textbf{Toepassing} & \textbf{Theorie} \\
        \midrule
        \endfirsthead
                
        \multicolumn{4}{c}{\tablename\ \thetable\ -- vervolg van vorige pagina} \\
        \toprule
        \multicolumn{2}{l}{\textbf{Principe}} & \textbf{Toepassing} & \textbf{Theorie} \\
        \midrule
        \endhead
                
        \midrule \multicolumn{4}{r}{Vervolg op volgende pagina} \\
        \endfoot
                
        \bottomrule
        \caption[Overzicht principes]{Overzicht van development gerelateerde ontwerpprincipes voor stressreducerend en -voorkomend ontwerp binnen de ontwikkeling van apps.}
        \label{tab:principes-dev}
        \endlastfoot
        \addlinespace[6pt]
        
        %==============================
        1.
        &
        Privacy by Design 
        & 
        \textit{"Wordt de privacy van de gebruiker beschermd zonder dat deze extra acties moet ondernemen?"}
        
        \begin{itemize}
            \item Verzamel enkel noodzakelijke gegevens
            \item Voorzie transparante privacy-instellingen
            \item Beveilig gevoelige data standaard
        \end{itemize} 
        & 
        {\scriptsize
            \ref{subsec:potentieel-en-randvoorwaarden} Privacy als randvoorwaarde voor m-health apps\par\bigskip
            \ref{subsec:kritieken-apps} Kritieken op apps: datalekken en gebrek aan beveiliging
        }
        \\
        
        \addlinespace[20pt]
        
        2.
        &
        Snel en eenvoudig distribueerbaar
        & 
        \textit{"Is je app snel vindbaar en makkelijk downloadbaar?"}
        
        \begin{itemize}
            \item Kies een eenvoudig te vinden distributieplatform
            \item Beperk de bestandsgrootte
            \item Gebruik de juiste technologie in functie van distributie, rekening houdend met verschillen in vindbaarheid en een installatiedrempel
        \end{itemize}
        & 
        {\scriptsize
            \ref{subsec:engagement-retentie} Engagement- en retentieproblematiek: drempels bij installatie\par\bigskip
            \ref{sec:technologieselectie} Technologie-selectie
        } 
        \\
        
        3.
        &
        Ondersteunende foutafhandeling
        & 
        \textit{"Helpt deze foutmelding de gebruiker verder zonder frustratie of onzekerheid te creëren?"}
        
        \begin{itemize}
            \item Begrijpelijke, niet-technische foutmeldingen
            \item Geen schuldinducerende formuleringen
            \item Bied altijd een duidelijke volgende stap
            \item Geen crashes of blokkades
            \item Fallback bij netwerk- of systeemproblemen
            \item Snelle en consistente laadtijden
        \end{itemize}
        & 
        {\scriptsize
            \ref{subsubsec:psychbehoeften-z} / \ref{subsec:uiux-algemeen} SDT: competentie\par\bigskip
            \ref{subsec:kritieken-apps} Kritieken op apps: frustratie door slechte UX en technische instabiliteit
        } 
        \\
        
        \addlinespace[20pt]
        
        4.
        &
        Inclusieve toegankelijkheid
        & 
        \textit{"Kan elke gebruiker zelfstandig gebruikmaken van de app?"}
        
        \begin{itemize}
            \item Volg de WCAG-richtlijnen
        \end{itemize}
        & 
        {\scriptsize
            \ref{subsubsec:psychbehoeften-z} / \ref{subsec:uiux-algemeen} SDT: autonomie\par\bigskip
            \ref{subsec:kritieken-apps} Kritieken op apps: slechte usability
        } 
        \\
        
        \addlinespace[20pt]
        
        5.
        &
        Beperk digitale interrupties
        & 
        \textit{"Respecteert de app rustmomenten van de gebruiker?"}
        
        \begin{itemize}
            \item Notificaties optioneel en instelbaar
            \item Geen vaste pushmeldingen zonder toestemming
            \item Lage en aanpasbare frequentie
            \item Geen dwingende of urgente formuleringen
        \end{itemize}
        & 
        {\scriptsize
            \ref{subsubsec:psychbehoeften-z} / \ref{subsec:uiux-algemeen} SDT: autonomie\par\bigskip
            \ref{subsec:info-overload} \acrshort{io}\par
            \ref{subsec:technostress} Technostress\par
            \ref{subsec:online-vigilance} Online vigilance\par\bigskip
            \ref{subsec:notificatie-strat} Notificatie strategieën
        } 
        \\
        
        \addlinespace[20pt]
        
    \end{longtable}
}

\newpage

\subsubsection{Design (UI/UX)}
\label{subsubsec:principes-design}

{
    \small
    \setlength\tabcolsep{3pt}
    \begin{longtable}{
            >{\raggedright\arraybackslash}p{0.2cm}  % Nummer
            >{\raggedright\arraybackslash}p{3.3cm}  % Principe (korte naam regel)
            @{\hspace{0.5cm}}
            >{\raggedright\arraybackslash}p{8.4cm}  % Toepassing (hoe toe te passen)
            @{\hspace{0.5cm}}
            >{\raggedright\arraybackslash}p{2.7cm}  % Theoretische basis (waarop werd dit gebaseerd?)
        }
        
        \toprule
        \multicolumn{2}{l}{\textbf{Principe}} & \textbf{Toepassing} & \textbf{Theorie} \\
        \midrule
        \endfirsthead
        
        \multicolumn{4}{c}{\tablename\ \thetable\ -- vervolg van vorige pagina} \\
        \toprule
        \multicolumn{2}{l}{\textbf{Principe}} & \textbf{Toepassing} & \textbf{Theorie} \\
        \midrule
        \endhead
        
        \midrule \multicolumn{4}{r}{Vervolg op volgende pagina} \\
        \endfoot
        
        \bottomrule
        \caption[Overzicht principes]{Overzicht van ontwerpgerelateerde ontwerpprincipes voor stressreducerend en -voorkomend ontwerp binnen de ontwikkeling van apps.}
        \label{tab:principes-design}
        \endlastfoot
        \addlinespace[6pt]
        
        %==============================
        6.
        &
        Behoud overzicht
        & 
        \textit{"Is het voor de gebruiker onmiddellijk duidelijk wat de belangrijkste informatie is en wat het doel van de pagina is?"}
        
        \begin{itemize}
            \item Voorzie een logische visuele hiërarchie
            \item Hanteer een consistente lay-out over alle schermen
            \item Beperk aantal elementen per scherm
            \item Highlight de belangrijkste actie
            \item Vermijd een ingewikkelde navigatie
        \end{itemize} 
        & 
        {\scriptsize
            \ref{subsubsec:psychbehoeften-z} / \ref{subsec:uiux-algemeen} SDT: competentie\par\bigskip
            \ref{subsec:info-overload} \acrshort{io}\par\bigskip
            \ref{subsec:kritieken-apps} Kritieken op apps: complexe navigatie en overvolle interfaces
        }
        \\
        
        \addlinespace[20pt]
        
        7.
        &
        Minimaliseer cognitieve belasting 
        & 
        \textit{"Moet de gebruiker te veel informatie tegelijk verwerken of onthouden?"}
        
        \begin{itemize}
            \item Informatie opgesplitst in kleine stukken (chunking)
            \item Pas progressive disclosure toe
            \item Voorzie korte, scanbare teksten
            \item Zorg voor 1 primaire taak per scherm
        \end{itemize} 
        & 
        {\scriptsize
            \ref{subsubsec:psychbehoeften-z} / \ref{subsec:uiux-algemeen} SDT: competentie\par\bigskip
            \ref{subsec:clt} Cognitive Load Theory\par\bigskip
            \ref{subsec:info-overload} \acrshort{io}\par\bigskip
            \ref{subsec:uiux-algemeen} UI/UX en cognitieve belasting
        }
        \\
        
        \addlinespace[20pt]
        
        8.
        &
        Vermijd visuele overprikkeling
        & 
        \textit{"Draagt het visuele ontwerp bij aan mentale rust?"}
        
        \begin{itemize}
            \item Rustig kleurenpalet (zachte kleuren zoals: blauw, groen, wit)
            \item Geen harde signaalkleuren (rood, oranje) zonder doel
            \item Maak gebruik van maximum 5 basiskleuren
            \item Maak gebruik van maximum 2 verschillende fonts
            \item Beperk animaties
            \item Gebruik afbeeldingen en video alleen ter ondersteuning van informatie
        \end{itemize} 
        & 
        {\scriptsize
            \ref{subsubsec:psychbehoeften-z} / \ref{subsec:uiux-algemeen} SDT: competentie\par\bigskip
            \ref{subsec:info-overload} \acrshort{io}\par\bigskip
            \ref{subsec:kritieken-apps} Kritieken op apps: overvolle interfaces en felle visuele elementen\par\bigskip
            \ref{subsec:uiux-algemeen} UI/UX en cognitieve belasting
        }
        \\
        
        \addlinespace[20pt]
        
        9.
        &
        Techno-simplicity 
        & 
        \textit{"Kan de gebruiker de applicatie intuïtief, snel en zonder extra hulp gebruiken?"}
        
        \begin{itemize}
            \item Voorzie een korte, optionele onboarding
            \item Hanteer consistente knoppen, iconen en terminologie
            \item Maak gebruik van tooltips of een inline uitleg bij onduidelijkere elementen
        \end{itemize} 
        & 
        {\scriptsize
            \ref{subsubsec:psychbehoeften-z} / \ref{subsec:uiux-algemeen} SDT: competentie\par\bigskip
            \ref{subsec:technostress} Technostress: techno-complexity\par\bigskip
            \ref{subsec:kritieken-apps} Kritieken op apps: complexe navigatie en slechte usability\par\bigskip
            \ref{subsec:engagement-retentie} Engagement- en retentieproblematiek: factoren voor adoptie
        }
        \\
        
        \addlinespace[20pt]
    \end{longtable}
}

\vspace{1em}

\subsubsection{Interaction}
\label{subsubsec:principes-interaction}

{
    \small
    \setlength\tabcolsep{3pt}
    \begin{longtable}{
            >{\raggedright\arraybackslash}p{0.3cm}  % Nummer
            >{\raggedright\arraybackslash}p{3.2cm}  % Principe (korte naam regel)
            @{\hspace{0.5cm}}
            >{\raggedright\arraybackslash}p{8.4cm}  % Toepassing (hoe toe te passen)
            @{\hspace{0.5cm}}
            >{\raggedright\arraybackslash}p{2.7cm}  % Theoretische basis (waarop werd dit gebaseerd?)
        }
        
        \toprule
        \multicolumn{2}{l}{\textbf{Principe}} & \textbf{Toepassing} & \textbf{Theorie} \\
        \midrule
        \endfirsthead
        
        \multicolumn{4}{c}{\tablename\ \thetable\ -- vervolg van vorige pagina} \\
        \toprule
        \multicolumn{2}{l}{\textbf{Principe}} & \textbf{Toepassing} & \textbf{Theorie} \\
        \midrule
        \endhead
        
        \midrule \multicolumn{4}{r}{Vervolg op volgende pagina} \\
        \endfoot
        
        \bottomrule
        \caption[Overzicht principes]{Overzicht van interactiegerelateerde ontwerpprincipes voor stressreducerend en -voorkomend ontwerp binnen de ontwikkeling van apps.}
        \label{tab:principes-interaction}
        \endlastfoot
        \addlinespace[6pt]
        
        %==============================
        10.
        &
        Stimuleer gezond gebruik
        & 
        \textit{"Kan de gebruiker de app makkelijk wegleggen?"}
        
        \begin{itemize}
            \item Vermijd eindeloze feeds
            \item Vermijd mechanismen die terugkeer afdwingen (vb. streaks)
            \item Beperk notificaties
        \end{itemize} 
        & 
        {\scriptsize
            \ref{subsubsec:psychbehoeften-z} / \ref{subsec:uiux-algemeen} SDT: autonomie\par\bigskip
            \ref{subsec:online-vigilance} Online vigilance\par\bigskip
            \ref{subsec:engagement-retentie} Engagement- en retentieproblematiek: strategieën voor een hogere retentie\par\bigskip
            \ref{subsec:notificatie-strat} Notificatie-strategieën
        }
        \\
        
        \addlinespace[20pt]
        
        11.
        &
        Zachte gedragssturing
        & 
        \textit{"Ondersteunt de applicatie gewenst gedrag zonder de gebruiker onder druk te zetten?"}
        
        \begin{itemize}
            \item Voorzie contextafhankelijke herinneringen
            \item Geef de mogelijkheid tot zelfmonitoring
            \item Geef de mogelijkheid tot uitstel/overslaan
            \item Stem ondersteuning af op voorkeuren/gebruikerscontext
        \end{itemize} 
        & 
        {\scriptsize
            \ref{subsubsec:psychbehoeften-z} / \ref{subsec:uiux-algemeen} SDT: autonomie\par\bigskip
            \ref{subsec:engagement-retentie} Engagement- en retentieproblematiek: strategieën voor een hogere retentie\par\bigskip
            \ref{subsec:motivatie-betrokkenheid-strat} Motivatie- en betrokkenheidsstrategieën\par\bigskip
            \ref{subsec:notificatie-strat} Notificatie-strategieën
        }
        \\
        
        \addlinespace[20pt]
        
        12.
        &
        Ondersteun autonomie 
        & 
        \textit{"Kan de gebruiker zelf keuzes maken en controle uitoefenen over zijn gebruik van de applicatie?"}
        
        \begin{itemize}
            \item Bied betekenisvolle keuzes aan
            \item Bied instelbare voorkeuren aan
            \item Bied personalisatie-opties aan
            \item Gebruik uitnodigende en niet-dwingende taal (vermijd “moeten”)
        \end{itemize} 
        & 
        {\scriptsize
            \ref{subsubsec:psychbehoeften-z} / \ref{subsec:uiux-algemeen} SDT: autonomie\par\bigskip
            \ref{subsec:kenmerken-generatieZ} Kenmerken van generatie Z: persoonlijkheid\par\bigskip
            \ref{subsec:engagement-retentie} Engagement- en retentieproblematiek: factoren voor adoptie\par\bigskip
            \ref{subsec:motivatie-betrokkenheid-strat} Motivatie- en betrokkenheidsstrategieën
        }
        \\
        
        \addlinespace[20pt]
        
        13.
        &
        Ondersteun competentiegevoel 
        & 
        \textit{"Ervaart de gebruiker voldoende succeservaringen en begrijpt hij het effect van zijn acties?"}
        
        \begin{itemize}
            \item Deel taken op in kleine stappen
            \item Geef directe feedback na acties
            \item Maak vooruitgang zichtbaar
            \item Stem, indien van toepassing, de moeilijkheid af op het niveau van de gebruiker
        \end{itemize} 
        & 
        {\scriptsize
            \ref{subsubsec:psychbehoeften-z} / \ref{subsec:uiux-algemeen} SDT: competentie\par\bigskip
            \ref{subsec:kenmerken-generatieZ} Kenmerken van generatie Z: persoonlijkheid\par\bigskip
            \ref{subsec:motivatie-betrokkenheid-strat} Motivatie- en betrokkenheidsstrategieën: bidirectional engagement
        }
        \\
        
        \addlinespace[20pt]
        
        14.
        &
        Stimuleer authentieke verbondenheid
        & 
        \textit{"Kan de gebruiker op een veilige en ondersteunende manier ervaringen delen zonder sociale druk?"}
        
        \begin{itemize}
            \item Geef de mogelijkheid om ervaringen te delen
            \item Stel duidelijke communityrichtlijnen op
            \item Vermijd sociale vergelijking of competitie
            \item Voorzie mogelijkheid tot contact met een echte persoon (bv. hulpverlener, peer support, …)
            \item Gebruik AI of chatbots enkel als aanvullende ondersteuning, niet als vervanging van menselijk contact
        \end{itemize} 
        & 
        {\scriptsize
            \ref{subsubsec:psychbehoeften-z} / \ref{subsec:uiux-algemeen} SDT: verbondenheid\par\bigskip
            \ref{subsec:kenmerken-generatieZ} Kenmerken van generatie Z: persoonlijkheid\par\bigskip
            \ref{subsec:risicofactoren-twintigers} Specifieke risicofactoren: maatschappelijke druk en beschermende patronen\par\bigskip
            \ref{subsec:kritieken-apps} Kritieken op apps: gebrek aan veilige en persoonlijke interactie\par\bigskip
            \ref{subsec:engagement-retentie} Engagement- en retentieproblematiek: strategie voor een hogere retentie
        }
        \\
        
        \addlinespace[20pt]
        
        15.
        &
        Betekenisvolle, prestatievrije motivatie
        & 
        \textit{"Ondersteunt de applicatie motivatie op een betekenisvolle manier zonder prestatiedruk of competitie?"}
        
        \begin{itemize}
            \item Voorzie beloningen gekoppeld aan echte vooruitgang
            \item Vermijd competitieve rankings
            \item Focus op persoonlijke groei en inspanning
            \item Gebruik narratieven of personages voor betrokkenheid
            \item Werk, indien van toepassing, met doelen of voortgang die geëvalueerd en bijgestuurd kunnen worden (\acrshort{smarter})
        \end{itemize} 
        & 
        {\scriptsize
            \ref{subsubsec:psychbehoeften-z} / \ref{subsec:uiux-algemeen} SDT: competentie en autonomie\par\bigskip
            \ref{subsec:risicofactoren-twintigers} Specifieke risicofactoren: maatschappelijke druk en beschermende patronen\par\bigskip
            \ref{subsec:engagement-retentie} Engagement- en retentieproblematiek: strategie voor een hogere retentie\par\bigskip
            \ref{subsec:motivatie-betrokkenheid-strat} Motivatie- en betrokkenheidsstrategieën
        }
        \\
        
    \end{longtable}
}

\vspace{1em}

\subsubsection{Content}
\label{subsubsec:principes-content}

{
    \small
    \setlength\tabcolsep{3pt}
    \begin{longtable}{
            >{\raggedright\arraybackslash}p{0.3cm}  % Nummer
            >{\raggedright\arraybackslash}p{3.2cm}  % Principe (korte naam regel)
            @{\hspace{0.5cm}}
            >{\raggedright\arraybackslash}p{8.4cm}  % Toepassing (hoe toe te passen)
            @{\hspace{0.5cm}}
            >{\raggedright\arraybackslash}p{2.7cm}  % Theoretische basis (waarop werd dit gebaseerd?)
        }
        
        \toprule
        \multicolumn{2}{l}{\textbf{Principe}} & \textbf{Toepassing} & \textbf{Theorie} \\
        \midrule
        \endfirsthead
        
        \multicolumn{4}{c}{\tablename\ \thetable\ -- vervolg van vorige pagina} \\
        \toprule
        \multicolumn{2}{l}{\textbf{Principe}} & \textbf{Toepassing} & \textbf{Theorie} \\
        \midrule
        \endhead
        
        \midrule \multicolumn{4}{r}{Vervolg op volgende pagina} \\
        \endfoot
        
        \bottomrule
        \caption[Overzicht principes]{Overzicht van contentgerelateerde ontwerpprincipes voor stressreducerend en -voorkomend ontwerp binnen de ontwikkeling van apps.}
        \label{tab:principes-content}
        \endlastfoot
        \addlinespace[6pt]
        
        %==============================
        16.
        &
        Persoonlijke relevantie
        & 
        \textit{"Sluit de inhoud aan bij de behoeften en context van de gebruiker?"}
        
        \begin{itemize}
            \item Content aangepast aan gebruiker, voorkeuren en gedrag 
            \item Vermijd veel generieke informatie
        \end{itemize} 
        & 
        {\scriptsize
            \ref{subsubsec:psychbehoeften-z} / \ref{subsec:uiux-algemeen} SDT: autonomie\par\bigskip
            \ref{subsec:kritieken-apps} Kritieken op apps: gebrek aan personalisatie\par\bigskip
            \ref{subsec:engagement-retentie} Engagement- en retentieproblematiek: factoren voor adoptie\par\bigskip
            \ref{subsec:motivatie-betrokkenheid-strat} Motivatie- en betrokkenheidsstrategieën
        }
        \\
        
        \addlinespace[20pt]
        
        17.
        &
        Zichtbare meerwaarde
        & 
        \textit{"Wordt de meerwaarde van de applicatie snel duidelijk?"}
        
        \begin{itemize}
            \item Benoem bij onboarding het doel en de meerwaarde expliciet
            \item Herhaal de meerwaarde impliciet in de inhoud (bv. uitleg, feedback)
        \end{itemize} 
        & 
        {\scriptsize
            \ref{subsec:kritieken-apps} Kritieken op apps: onduidelijke meerwaarde\par\bigskip
            \ref{subsec:engagement-retentie} Engagement- en retentieproblematiek: factoren voor adoptie\par\bigskip
            \ref{subsec:motivatie-betrokkenheid-strat} Motivatie- en betrokkenheidsstrategieën
        }
        \\
        
        \addlinespace[20pt]
        
        18.
        &
        Duidelijke positionering
        & 
        \textit{"Is het voor de gebruiker duidelijk wat de applicatie wel en niet biedt?"}
        
        \begin{itemize}
            \item Beschrijf duidelijk de functionaliteiten en beperkingen van de app (bv. onboarding of app store)
            \item Vermijd vage of misleidende claims
            \item Gebruik duidelijke naamgeving en beschrijvingen van features
            \item Maak duidelijk voor welke doelgroep de applicatie bedoeld is
        \end{itemize} 
        & 
        {\scriptsize
            \ref{subsubsec:psychbehoeften-z} / \ref{subsec:uiux-algemeen} SDT: autonomie\par\bigskip
            \ref{subsec:kritieken-apps} Kritieken op apps: kloof tussen commerciële claims en effectieve werking\par\bigskip
            \ref{subsec:engagement-retentie} Engagement- en retentieproblematiek: factoren voor adoptie
        }
        \\
        
        \addlinespace[20pt]
        
        19.
        &
        Ondersteunende communicatie
        & 
        \textit{"Is de communicatie empathisch en niet-oordelend?"}
        
        \begin{itemize}
            \item Gebruik empathische, begrijpelijke taal
            \item Vermijd schuldinducerende boodschappen
            \item Vermijd normatieve uitspraken over geluk of succes
        \end{itemize} 
        & 
        {\scriptsize
            \ref{subsec:risicofactoren-twintigers} Specifieke risicofactoren: maatschappelijke druk en beschermende patronen\par\bigskip
            \ref{subsec:uiux-algemeen} UI/UX en cognitieve belasting
        }
        \\
        
        \addlinespace[20pt]
        
        
    \end{longtable}
}